\loeskopf{Einheit 2 von 5 · Satz vom Nullprodukt (Nr. 10--16)}

\erg{10}{a) ja \quad b) nein \quad c) ja \quad d) nein \quad e) ja}

\erg{11}{a) $x_1 = 2$, $x_2 = 9$ \quad b) $x_1 = 4$, $x_2 = 11$ \quad
c) $x_1 = 1$, $x_2 = 6$ \quad d) $x_1 = 5$, $x_2 = 12$ \quad
e) $x_1 = -7$, $x_2 = 3$ \quad f) $x_1 = 0$, $x_2 = -4$ \quad
g) $x = 6$ (beide Faktoren gleich, nur eine Lösung) \quad
h) $x \cdot (x + 9) = 0$; $x_1 = 0$, $x_2 = -9$ \quad
i) $x \cdot (x - 13) = 0$; $x_1 = 0$, $x_2 = 13$ \quad
j) $4x \cdot (x + 5) = 0$; $x_1 = 0$, $x_2 = -5$ \quad
k) $3x \cdot (x - 4) = 0$; $x_1 = 0$, $x_2 = 4$}

\erg{12}{a) $x^2 + 6x - 27 = 0$ \quad b) $x^2 + 3x - 10 = 18$, also
$x^2 + 3x - 28 = 0$ \quad c) $x^2 - 9x + 14 = 0$}

\erg{13}{a) wahr ($(-6) \cdot 4 = -24$) \quad b) wahr ($(-4) \cdot 5 = -20$) \quad
c) falsch ($3 \cdot (-3) = -9 \neq 9$)}

\erg{14}{Fehler: Sara hat durch $x$ geteilt. Das ist nur erlaubt, wenn $x$ nicht null
ist -- dabei geht die Lösung $x = 0$ verloren. Richtig: alles auf eine Seite,
$x^2 - 11x = 0$, ausklammern $x \cdot (x - 11) = 0$, also $x_1 = 0$, $x_2 = 11$.
\quad a) $x^2 - 6x = 0$; $x \cdot (x - 6) = 0$; $x_1 = 0$, $x_2 = 6$}

\erg{15}{a) Null mal irgendeine Zahl ergibt immer null; ist ein Faktor null, ist das
ganze Produkt null. \quad b) Beim Teilen durch $x$ nimmt sie an, dass $x$ nicht null
ist. Gerade $x = 0$ ist aber eine Lösung; sie fällt beim Teilen weg. Übrig bleibt nur
$x = 8$. \quad c) Die erste hat zwei Lösungen, weil die beiden Faktoren für
verschiedene $x$ null werden; die zweite hat nur eine, weil beide Faktoren gleich sind
und nur für $x = -3$ null werden.}

\erg{16}{a) $x = 0$ und $x = 6$ \quad b) $x = -3$, denn $(-3) \cdot 4 = -12$}
